\documentclass[12pt,a4paper]{article}

\usepackage[T2A]{fontenc}
\usepackage[utf8]{inputenc}
\usepackage[russian]{babel}
\usepackage{amsmath}
\usepackage{amssymb}
\usepackage{array}
\usepackage{booktabs}
\usepackage{caption}
\usepackage{enumitem}
\usepackage{float}
\usepackage{geometry}
\usepackage{graphicx}
\usepackage{hyperref}
\usepackage{longtable}
\usepackage{multirow}
\usepackage{tabularx}

\geometry{margin=2.2cm}
\hypersetup{hidelinks}
\captionsetup{font=normalsize}
\graphicspath{{report_assets/}}
\setlist{nosep}

\newcolumntype{Y}{>{\centering\arraybackslash}X}

\title{Отчёт по заданию 3\\Построение нейронных сетей для прогноза цены и точек разворота}
\author{}
\date{}

\begin{document}
\maketitle
\tableofcontents
\newpage

\section{Постановка задачи}
Требуется построить три прогностические модели по минутно-секундным данным торгов акциями Сбербанка:
\begin{enumerate}
    \item модель для прогноза обращённого медленного стохастического осциллятора \texttt{Slow \%K};
    \item модель для прогноза нижних точек разворота;
    \item модель для прогноза верхних точек разворота.
\end{enumerate}

После обучения моделей нужно:
\begin{enumerate}
    \item сравнить несколько трёхслойных и четырёхслойных архитектур;
    \item выбрать лучшие сети по качеству на валидации;
    \item подобрать торговые пороги;
    \item проверить торговые результаты на новом дне и отдельно оценить устойчивость на стресс-тесте.
\end{enumerate}

\section{План работы}
Работа выполнена по следующему плану:
\begin{enumerate}
    \item из потока сделок построить секундные бары и подготовить признаки;
    \item задать целевые переменные для цены и точек разворота;
    \item обучить несколько архитектур MLP-сетей;
    \item выбрать лучшие модели по валидационному качеству;
    \item подобрать пороги торговых сигналов;
    \item проверить модели на новом торговом дне;
    \item отдельно оценить поведение стратегий на дне с иным рыночным режимом;
    \item сформулировать содержательные выводы по качеству и устойчивости моделей.
\end{enumerate}

\section{Данные и схема эксперимента}
\subsection{Подготовка данных}
Из потока сделок построены секундные бары с полями \texttt{open}, \texttt{high}, \texttt{low}, \texttt{close}, \texttt{volume}, \texttt{amount}, количеством сделок и знаковым объёмом. Пропущенные секунды по цене заполнялись предыдущим значением, а объёмные характеристики занулялись.

Для анализа использованы три торговых дня:
\begin{itemize}
    \item \textbf{14.03.2013} --- день обучения и валидации;
    \item \textbf{15.03.2013} --- основной внешний тест;
    \item \textbf{16.01.2013} --- дополнительный стресс-тест.
\end{itemize}

Внутри дня 14.03.2013 первая половина сессии использована для обучения, вторая --- для валидации. Такая схема позволяет сначала выбирать архитектуры и пороги на одном дне, а затем проверять результат на следующем торговом дне.

\begin{figure}[H]
    \centering
    \includegraphics[width=\textwidth]{overview_days.png}
    \caption{Общий обзор данных: цена, нормированная траектория, ликвидность и локальная амплитуда движения}
\end{figure}

Из рисунка видно, что мартовские дни близки по масштабу движения, тогда как январский день заметно спокойнее и менее ликвиден. Это делает его удобным стресс-тестом на смену рыночного режима.

\subsection{Признаки и целевые переменные}
Для модели цены использовался вектор
\[
x_k = \frac{x_{i-n_k} - x_i}{\sqrt{n_k}},
\]
где
\[
n_k \in \{10, 20, 30, 40, 50, 60, 80, 100, 120, 160, 200, 240, 320, 340, 400, 480, 640, 960\},
\]
после чего весь вектор нормировался по евклидовой норме.

Для моделей точек разворота к этому набору были добавлены простые динамические признаки: лаговые изменения цены, отклонения от EMA, локальная волатильность, усреднённая активность и дисбаланс объёма.

Целевая переменная для регрессии --- обращённый медленный осциллятор
\[
\text{InvSlowK} = 100 - \text{Slow \%K}.
\]

Нижняя точка разворота задавалась меткой \(1\), если текущая цена ниже трёх предыдущих и десяти последующих значений. Верхняя точка разворота задавалась аналогично: текущая цена должна быть выше трёх предыдущих и десяти последующих значений.

\section{Построение моделей}
\subsection{Архитектуры}
Для регрессионной модели были проверены две трёхслойные и две четырёхслойные архитектуры:
\begin{itemize}
    \item \texttt{3-layer / 16};
    \item \texttt{3-layer / 32};
    \item \texttt{4-layer / 24-12};
    \item \texttt{4-layer / 32-16}.
\end{itemize}

Для моделей минимумов и максимумов были проверены архитектуры:
\begin{itemize}
    \item \texttt{3-layer / 24};
    \item \texttt{3-layer / 48};
    \item \texttt{4-layer / 32-16};
    \item \texttt{4-layer / 48-24}.
\end{itemize}

Во всех случаях использовались сети прямого распространения из библиотеки scikit-learn. Для цены решалась задача регрессии, для минимумов и максимумов --- задача бинарной классификации.

\subsection{Критерии отбора}
Для модели цены основной критерий --- корреляция на валидации. Дополнительно контролировались скорректированная корреляция и средняя абсолютная ошибка.

Для моделей минимумов и максимумов основной критерий --- \texttt{ROC-AUC}. Это разумный выбор для редких событий: важно не только абсолютное значение вероятности, но и способность модели правильно ранжировать наблюдения по силе сигнала.

\section{Результаты обучения}
\subsection{Модель цены}
\begin{table}[H]
    \centering
    \small
    \caption{Сравнение архитектур для прогноза обращённого медленного \%K}
    \resizebox{\textwidth}{!}{
    \begin{tabular}{lccccccc}
        \toprule
        Название сети & Слои & Размер & Связи & Corr train & Corr valid & Corr\(_c\) valid & MAE valid \\
        \midrule
        3-layer / 32 & 3-layer & 32 & 608 & 0.8231 & 0.7258 & 0.7120 & 21.3926 \\
        3-layer / 16 & 3-layer & 16 & 304 & 0.8132 & 0.7247 & 0.7179 & 21.3899 \\
        4-layer / 24-12 & 4-layer & 24-12 & 732 & 0.8496 & 0.6885 & 0.6688 & 22.4232 \\
        4-layer / 32-16 & 4-layer & 32-16 & 1104 & 0.8655 & 0.6844 & 0.6530 & 22.4196 \\
        \bottomrule
    \end{tabular}}
\end{table}

Лучшая трёхслойная сеть --- \texttt{3-layer / 32}. Лучшая четырёхслойная сеть --- \texttt{4-layer / 24-12}. При этом уже на таблице 1 видно, что увеличение глубины не дало выигрыша по качеству на валидации: более простая архитектура оказалась устойчивее.

\subsection{Модель минимумов}
\begin{table}[H]
    \centering
    \small
    \caption{Сравнение архитектур для прогноза нижних точек разворота}
    \resizebox{\textwidth}{!}{
    \begin{tabular}{lccccccc}
        \toprule
        Название сети & Слои & Размер & Связи & Corr valid & Corr\(_c\) valid & ROC-AUC valid & Положительных \\
        \midrule
        3-layer / 48 & 3-layer & 48 & 1824 & 0.2086 & 0.0000 & 0.9232 & 94 \\
        3-layer / 24 & 3-layer & 24 & 912 & 0.1903 & 0.0000 & 0.9087 & 94 \\
        4-layer / 48-24 & 4-layer & 48-24 & 2952 & 0.1944 & 0.0000 & 0.8965 & 94 \\
        4-layer / 32-16 & 4-layer & 32-16 & 1712 & 0.1849 & 0.0000 & 0.8907 & 94 \\
        \bottomrule
    \end{tabular}}
\end{table}

Лучшая сеть для минимумов --- \texttt{3-layer / 48}. Здесь особенно хорошо видно, что линейная корреляция малоинформативна: событие редкое, поэтому главный критерий --- именно качество ранжирования по \texttt{ROC-AUC}.

\subsection{Модель максимумов}
\begin{table}[H]
    \centering
    \small
    \caption{Сравнение архитектур для прогноза верхних точек разворота}
    \resizebox{\textwidth}{!}{
    \begin{tabular}{lccccccc}
        \toprule
        Название сети & Слои & Размер & Связи & Corr valid & Corr\(_c\) valid & ROC-AUC valid & Положительных \\
        \midrule
        3-layer / 24 & 3-layer & 24 & 912 & 0.2555 & 0.0630 & 0.9382 & 134 \\
        3-layer / 48 & 3-layer & 48 & 1824 & 0.2592 & 0.0000 & 0.9294 & 134 \\
        4-layer / 32-16 & 4-layer & 32-16 & 1712 & 0.2262 & 0.0000 & 0.9008 & 134 \\
        4-layer / 48-24 & 4-layer & 48-24 & 2952 & 0.2298 & 0.0000 & 0.8813 & 134 \\
        \bottomrule
    \end{tabular}}
\end{table}

Лучшая сеть для максимумов --- \texttt{3-layer / 24}. Как и в модели минимумов, более глубокие архитектуры здесь не дают улучшения.

\section{Качество моделей на валидации, тесте и стресс-тесте}
\subsection{Модель цены}
\begin{figure}[H]
    \centering
    \includegraphics[width=\textwidth]{price_distribution.png}
    \caption{Диагностика распределения регрессионной модели: распределение факта и прогноза, калибровка и ошибки}
\end{figure}

\begin{figure}[H]
    \centering
    \includegraphics[width=\textwidth]{price_timeseries.png}
    \caption{Временная картина регрессионной модели: весь день, фокусное окно и увеличенный фрагмент}
\end{figure}

Ключевые результаты лучшей регрессионной сети \texttt{3-layer / 32}:
\begin{itemize}
    \item на валидации: корреляция \(0.7258\), MAE \(21.39\), RMSE \(27.21\);
    \item на новом дне: корреляция \(0.7008\), MAE \(22.69\), RMSE \(28.90\);
    \item на стресс-тесте: корреляция \(0.7592\), MAE \(21.84\), RMSE \(27.52\).
\end{itemize}

Регрессионная сеть показывает достаточно стабильное качество и не деградирует при переносе на другой день. Это сильный результат: модель не просто подгоняет конкретный фрагмент сессии, а удерживает разумную структуру прогноза и на новом дне, и на более спокойном январском рынке.

\subsection{Модели точек разворота}
\begin{figure}[H]
    \centering
    \includegraphics[width=\textwidth]{turn_quality.png}
    \caption{Качество моделей минимумов и максимумов: ROC-кривые и кривые захвата событий}
\end{figure}

\begin{figure}[H]
    \centering
    \includegraphics[width=\textwidth]{turn_signals_price.png}
    \caption{Сигналы модели точек разворота по цене: весь день, фокусное окно и zoom}
\end{figure}

\begin{figure}[H]
    \centering
    \includegraphics[width=\textwidth]{turn_signals_prob.png}
    \caption{Вероятности модели точек разворота: весь день, фокусное окно и zoom}
\end{figure}

Результаты лучших классификаторов:
\begin{itemize}
    \item минимумы, \texttt{3-layer / 48}: \texttt{ROC-AUC} \(0.9232\) на валидации, \(0.9058\) на новом дне и \(0.9116\) на стресс-тесте;
    \item максимумы, \texttt{3-layer / 24}: \texttt{ROC-AUC} \(0.9382\) на валидации, \(0.9420\) на новом дне и \(0.9540\) на стресс-тесте.
\end{itemize}

При этом \texttt{Average Precision} заметно ниже: для минимумов на новом дне она равна \(0.0764\), для максимумов --- \(0.1494\). Это закономерно: точки разворота редки, поэтому модель хорошо ранжирует события, но детектор точного момента остаётся шумным. Иначе говоря, сеть лучше подходит для отбора сильных кандидатов на вход, чем для точного бинарного решения в каждой секунде.

\section{Подбор порогов и торговые результаты}
Порог для модели цены подбирался в диапазоне от 50 до 90, для моделей точек разворота --- от 20 до 80 с шагом 2. Три типа входа в сделку:
\begin{itemize}
    \item \texttt{market};
    \item \texttt{limit};
    \item \texttt{stop}.
\end{itemize}

Порог подбирался на валидационном интервале и затем переносился без изменения на новый день и на стресс-тест.

\subsection{Применение на новом дне}
\begin{table}[H]
    \centering
    \scriptsize
    \caption{Применение модели цены на новом дне 15.03.2013}
    \resizebox{\textwidth}{!}{
    \begin{tabular}{lcccccc}
        \toprule
        Модель & Размер сети & Приказ & Порог & Доходность, \% годовых & Достоверность, \% & Прибыльных, \% / сделок \\
        \midrule
        Обращённый Slow \%K & 32 & limit & 50 & 148.03 & 77.49 & 43.88 / 139 \\
        Обращённый Slow \%K & 24-12 & limit & 67 & -14.64 & 46.87 & 41.67 / 132 \\
        Обращённый Slow \%K & 32 & market & 65 & -100.96 & 30.90 & 38.99 / 159 \\
        Обращённый Slow \%K & 24-12 & market & 73 & -13.06 & 47.47 & 39.62 / 159 \\
        Обращённый Slow \%K & 32 & stop & 51 & -155.41 & 20.00 & 35.04 / 137 \\
        Обращённый Slow \%K & 24-12 & stop & 56 & -409.11 & 1.57 & 31.29 / 147 \\
        \bottomrule
    \end{tabular}}
\end{table}

\begin{table}[H]
    \centering
    \scriptsize
    \caption{Применение модели точек разворота на новом дне 15.03.2013}
    \resizebox{\textwidth}{!}{
    \begin{tabular}{lcccccc}
        \toprule
        Модель & Размер сети & Приказ & Порог & Доходность, \% годовых & Достоверность, \% & Прибыльных, \% / сделок \\
        \midrule
        Точки разворота (min/max) & 48 + 24 & limit & 56 & 129.26 & 83.97 & 47.46 / 59 \\
        Точки разворота (min/max) & 48-24 + 32-16 & limit & 48 & 229.78 & 94.40 & 50.00 / 72 \\
        Точки разворота (min/max) & 48 + 24 & market & 34 & 72.29 & 64.07 & 40.71 / 140 \\
        Точки разворота (min/max) & 48-24 + 32-16 & market & 30 & 190.44 & 82.16 & 43.92 / 148 \\
        Точки разворота (min/max) & 48 + 24 & stop & 50 & -97.65 & 28.27 & 39.29 / 112 \\
        Точки разворота (min/max) & 48-24 + 32-16 & stop & 30 & 26.26 & 55.57 & 41.41 / 128 \\
        \bottomrule
    \end{tabular}}
\end{table}

\begin{figure}[H]
    \centering
    \includegraphics[width=\textwidth]{strategy_results.png}
    \caption{Сравнение лучших стратегий на новом дне}
\end{figure}

Лучший результат на новом дне дала стратегия по точкам разворота с комбинацией сетей \texttt{48-24 + 32-16} и типом входа \texttt{limit}. Её результат:
\[
229.78\% \text{ годовых}, \qquad 94.40\% \text{ статистической достоверности}, \qquad 72 \text{ сделки}.
\]

Сильный вывод из таблицы 4: входы \texttt{limit} систематически лучше \texttt{stop}, а для ряда моделей и лучше \texttt{market}. Это согласуется с природой задачи: модели пытаются ловить локальные экстремумы, поэтому более аккуратный вход вблизи предполагаемого разворота работает лучше агрессивного входа по факту движения.

\subsection{Стресс-тест на 16.01.2013}
На дополнительном стресс-тесте картина изменилась:
\begin{itemize}
    \item у модели цены лучшими оказались стратегии \texttt{limit}: \(107.18\%\) для сети \texttt{32} и \(159.76\%\) для сети \texttt{24-12};
    \item у моделей точек разворота устойчивость оказалась хуже: лучший результат дал только вариант \texttt{48 + 24 / limit} с \(37.71\%\), а остальные комбинации заметно просели, вплоть до сильно отрицательных доходностей.
\end{itemize}

Это важный практический результат. На соседнем мартовском дне модели разворота показывают максимальную доходность, но при смене режима именно регрессионная модель цены переносится заметно лучше.

\section{Обсуждение результатов}
Полученные результаты позволяют сделать несколько содержательных наблюдений.

\begin{enumerate}
    \item \textbf{Более глубокая сеть не дала преимущества.} Во всех трёх задачах лучшие результаты показали относительно простые архитектуры. Для данного объёма обучающей выборки это выглядит естественно: глубина добавляет гибкость, но не повышает качество обобщения.

    \item \textbf{Модель цены оказалась более устойчивой, чем модель точек разворота.} По ближайшему новому дню стратегия на разворотах действительно сильнее. Однако на более спокойном стрессовом дне регрессионная модель переносится заметно лучше. С практической точки зрения это аргумент в пользу более ``плавных'' целей, а не только редких экстремумов.

    \item \textbf{ROC-AUC высок, но этого недостаточно для торгового применения.} Для редких событий сети минимумов и максимумов хорошо ранжируют наблюдения, но это не означает, что точка входа автоматически будет точной и стабильной. Низкий \texttt{Average Precision} и ухудшение на стресс-тесте это подтверждают.

    \item \textbf{Тип приказа критически важен.} Один и тот же прогноз может давать совершенно разные торговые результаты при \texttt{market}, \texttt{limit} и \texttt{stop}. Поэтому оценивать нейросеть только по ML-метрикам недостаточно; её нужно проверять в торговом контуре.

    \item \textbf{Валидационный подбор порога даёт рабочий результат, но остаётся источником переобучения.} Пороги подбирались на одном интервале внутри дня. Это достаточно для учебной постановки, но в реальной задаче потребовало бы walk-forward схемы и более длинной истории.
\end{enumerate}

\section{Заключение}
В работе были построены три нейронные сети: регрессионная модель для обращённого медленного \%K и две классификационные модели для нижних и верхних точек разворота. Для каждой задачи были сравнены трёхслойные и четырёхслойные архитектуры, после чего выбраны лучшие сети по качеству на валидации.

Основные итоги работы:
\begin{itemize}
    \item для цены лучшей оказалась сеть \texttt{3-layer / 32};
    \item для минимумов --- \texttt{3-layer / 48};
    \item для максимумов --- \texttt{3-layer / 24};
    \item на новом дне максимальную доходность показала стратегия по точкам разворота \texttt{48-24 + 32-16 / limit} с результатом \(229.78\%\) годовых;
    \item на стресс-тесте более устойчивой оказалась регрессионная модель цены.
\end{itemize}

Главный содержательный вывод состоит в том, что для этой задачи нейросети действительно умеют извлекать полезный сигнал из секундных данных, но разные постановки оптимизируют разные свойства. Модель точек разворота лучше ловит локальные возможности на близком рыночном режиме, а модель цены даёт более устойчивое поведение при переносе на другой режим торговли.

\end{document}
